% section5_dynamic_model.tex

\section{Dynamic Model: Anticipatory Top-Down Pooling}
\label{sec:dynamic-model}

Sections~\ref{sec:model}--\ref{sec:dynamic} use a one-period anticipatory
exposure term. At a low-scrutiny history, choosing \(s_A\) today exposes the
sender to possible adverse repricing tomorrow. If the sender climbs down to
\(s_P\), tomorrow's penalty from today's ambitious signal is shed. Thus
withdrawal has dynamic value even though exposure is not a permanent stock.

The reduced-form high-type incentive constraint is
\[
D_H^L(q)=B_H-\beta q\mu_H\Phi_H.
\]
The dynamic extension in this section keeps that economics fixed and adds only
the missing law of motion. The transition probability \(q_t\) is a Markov state
observed by the sender. The paper's continuous-limit object is then not a new
formalism. Its potential is the stage-game stay advantage itself, and its
turning point is exactly the discrete IC threshold.

\subsection{The one-period Markov environment}
\label{subsec:one-period-markov-environment}

Time is discrete. The scrutiny regime is
\[
j_t\in\{L,H\},
\]
with \(L\) denoting ordinary evaluation and \(H\) denoting high scrutiny. In a
low-scrutiny period, the public Markov state
\[
q_t\in[0,1]
\]
is the conditional probability that tomorrow's regime is high scrutiny:
\[
q_t=\Pr(j_{t+1}=H\mid j_t=L,q_t).
\]
The process \(q_t\) evolves according to a common-knowledge Markov kernel
\[
Q(dq_{t+1}\mid q_t).
\]
The sender observes \((j_t,q_t)\) before choosing \(s_t\in\{s_A,s_P\}\).

\begin{assumption}[Common-knowledge Markov law]
\label{ass:common-knowledge-markov-law}
The kernel \(Q(dq_{t+1}\mid q_t)\), the regime-transition law, and the
sender's observation of \((j_t,q_t)\) are common knowledge. Conditional on the
public state, the law of future \((j,q)\)-paths does not depend on private
higher-order beliefs.
\end{assumption}

This assumption is the dynamic counterpart of the regime-conditional
continuity discipline in Section~\ref{sec:surprise}. If the law of future
\(q_t\)'s were private or itself belief-dependent, the sender would need
beliefs about the receiver's beliefs about future transition risk.
Assumption~\ref{ass:common-knowledge-markov-law} prevents that extra hierarchy
from entering the payoff technology.

\begin{lemma}[Intertemporal continuity]
\label{lem:intertemporal-continuity}
Under Assumptions~\ref{ass:conditional-continuity} and
\ref{ass:common-knowledge-markov-law}, the repeated one-period psychological
game is well posed. The dynamic payoff is continuous in the relevant
intertemporal belief hierarchy.
\end{lemma}

\begin{proof}
Fix a public history and a strategy profile. By
Lemma~\ref{lem:regime-continuity}, the one-period psychological payoff is
continuous in the relevant belief hierarchy within each fixed regime. The
action and type spaces are finite, and the public state \(q_t\) lies in a
compact interval. The one-period payoff is therefore bounded on each regime
slice. By Assumption~\ref{ass:common-knowledge-markov-law}, the conditional law
of future states is fixed and common knowledge. The repeated payoff is a
discounted expectation of bounded continuous one-period payoff functions under
a fixed measure. Discounting gives a uniformly convergent series, and uniform
limits of continuous functions are continuous. Hence the repeated
psychological game is well posed in the sense of dynamic psychological games
\citep{battigalli2009dynamic}.
\end{proof}

The model remains genuinely anticipatory. The sender chooses while
\(j_t=L\), but the payoff of \(s_A\) depends on the probability that the next
period will be \(H\). Climbing down is valuable because silence today removes
the adverse-repricing exposure that would otherwise be carried into tomorrow.

\subsection{The micro-derived resistance potential}
\label{subsec:micro-derived-potential}

The continuous-limit potential is not imposed. It is the high type's
one-period stay advantage:
\[
V_H(q):=D_H^L(q)=B_H-\beta q\mu_H\Phi_H.
\]
The high-type IC threshold is therefore
\[
q_H^*
=
\frac{B_H}{\beta\mu_H\Phi_H}.
\]
Equivalently,
\[
V_H(q_H^*)=0,
\qquad
V_H'(q)=-\beta\mu_H\Phi_H<0.
\]
Thus the zero of the continuum potential is exactly the discrete high-type IC
break:
\[
x_c=q_H^*.
\]

\begin{proposition}[The turning point is the IC threshold]
\label{prop:turning-point-is-ic}
In the one-period Markov environment, the continuum potential generated by the
high-type stage-game payoff has a unique simple zero at
\[
x_c=q_H^*.
\]
At this point the separating high-type incentive constraint switches sign.
\end{proposition}

\begin{proof}
The high type's stay advantage is
\[
V_H(q)=B_H-\beta q\mu_H\Phi_H.
\]
Since \(\beta>0\), \(\mu_H>0\), and \(\Phi_H>0\), the derivative is strictly
negative. Hence \(V_H\) has at most one zero. Solving \(V_H(q)=0\) gives
\[
q=\frac{B_H}{\beta\mu_H\Phi_H}=q_H^*.
\]
For \(q<q_H^*\), \(V_H(q)>0\), so the high type's ambitious signal is protected
by a positive stay advantage. For \(q>q_H^*\), \(V_H(q)<0\), so the high type's
IC for \(s_A\) fails. The potential's turning point is therefore the discrete
IC threshold.
\end{proof}

This proposition is the weld between the finite game and the continuous
Markov limit. If the continuum potential had an unrelated zero, the
Sturm--Liouville problem would describe a different model. Here the zero is
derived from the same payoff inequality that drives top-down pooling in
Sections~\ref{sec:model}--\ref{sec:dynamic}.

\subsection{Discrete withdrawal fields and the Airy limit}
\label{subsec:discrete-withdrawal-airy}

Let the interval of possible transition states be \(I=[a,b]\subseteq[0,1]\),
and let \(q_i=a+ih\) be a grid. A near-threshold high-type withdrawal field
\[
\psi_h=(\psi_{h,0},\ldots,\psi_{h,n+1})
\]
records the refined discrete resistance of high-type occupation of \(s_A\)
across nearby \(q\)-states. The field is not an additional preference
primitive. It is a smoothed continuum representation of the discrete stay
advantage \(V_H(q_i)\).

Define the discrete operator
\[
(H_h\psi_h)_i
=
-\sigma(\Delta_h\psi_h)_i+V_H(q_i)\psi_{h,i},
\qquad
\sigma>0,
\]
where \(\Delta_h\) is the second-difference operator. The quasi-static
refinement condition is that the field is near-null in the discrete dual norm:
\[
\|\psi_h\|_{L_h^2}=1,
\qquad
\|H_h\psi_h\|_{H_h^{-1}}\to 0.
\]
This is the finite-mesh analogue of a stationary Fokker--Planck balance for
the withdrawal density. The diffusion term \(\sigma\Delta_h\) captures local
dispersion in the \(q\)-state, while the potential term is the micro-derived
stay advantage.

\begin{proposition}[Continuous withdrawal equation]
\label{prop:continuous-withdrawal-equation}
Suppose the near-null condition above holds along a mesh refinement. Then,
after passing to a subsequence, the interpolated fields converge weakly to a
nonzero limit \(\Psi\) satisfying
\[
-\sigma\Psi''(q)+V_H(q)\Psi(q)=0
\]
in the weak Sturm--Liouville sense. Since \(V_H\) has a simple zero at
\(q_H^*\), the local normal form near \(q_H^*\) is Airy.
\end{proposition}

\begin{proof}
The argument is the standard discrete Sobolev compactness step for
sign-changing Sturm--Liouville forms. The near-null condition gives a uniform
discrete \(H^1\) bound after a Garding shift. Interpolating the fields and
using compactness yields a weak \(H^1\) and strong \(L^2\) limit \(\Psi\). The
dual residual vanishes against every smooth test function, so the limit
satisfies
\[
\sigma\int_I \Psi'(q)\eta'(q)\,dq
+
\int_I V_H(q)\Psi(q)\eta(q)\,dq
=0
\]
for every test function \(\eta\in H_0^1(I)\). This is the weak form of
\(-\sigma\Psi''+V_H\Psi=0\). By
Proposition~\ref{prop:turning-point-is-ic}, \(V_H(q_H^*)=0\) and
\[
V_H'(q_H^*)=-\beta\mu_H\Phi_H<0.
\]
Thus \(q_H^*\) is a simple turning point. With
\[
z=
\left(\frac{\beta\mu_H\Phi_H}{\sigma}\right)^{1/3}
(q-q_H^*),
\]
the local equation becomes, to first order,
\[
Y''(z)+zY(z)=0,
\]
whose local basis is
\[
Y(z)
=
C_{\operatorname{Ai}}\operatorname{Ai}(-z)
+
C_{\operatorname{Bi}}\operatorname{Bi}(-z).
\]
\end{proof}

The Airy profile is therefore not assumed. It appears only after the
stage-game potential has been derived and the discrete withdrawal field has
been passed to its continuum limit.

\subsection{Protected-side boundedness}
\label{subsec:protected-side-boundedness}

The local Airy basis alone does not select a branch. A final economic boundary
condition is needed.

\begin{assumption}[Protected-side boundedness]
\label{ass:protected-side-boundedness}
In the protected region \(q<q_H^*\), where the high type's stay advantage is
strictly positive, the refined resistance field remains uniformly bounded as
the mesh is refined. Economically, there is no exploding mass of near-withdrawal
pressure in states where the high type's ambitious IC is slack.
\end{assumption}

\begin{proposition}[Airy branch selection]
\label{prop:airy-branch-selection}
Under Assumption~\ref{ass:protected-side-boundedness}, the growing
\(\operatorname{Bi}\) component is ruled out on the protected side:
\[
C_{\operatorname{Bi}}=0.
\]
The local withdrawal-density profile near the threshold is therefore the
bounded Airy branch.
\end{proposition}

\begin{proof}
For \(q<q_H^*\), the Airy coordinate satisfies \(z<0\), so \(-z>0\). Along the
protected side, \(\operatorname{Bi}(-z)\) is the growing branch, while
\(\operatorname{Ai}(-z)\) is the bounded/decaying branch. If
\(C_{\operatorname{Bi}}\ne 0\), the resistance field would diverge as the
protected-side limit is approached. That violates
Assumption~\ref{ass:protected-side-boundedness}. Hence
\(C_{\operatorname{Bi}}=0\).
\end{proof}

Section~\ref{sec:dynamic-model} therefore performs the required extension
rather than adding a parallel formalism. One-period exposure gives the repeated
static IC
\[
B_H-\beta q\mu_H\Phi_H,
\]
the micro-derived potential is that same IC expression, the continuum turning
point is
\[
x_c=q_H^*,
\]
and the Airy limit follows from the weak Sturm--Liouville limit plus the
protected-side boundedness condition. The remaining step is not the existence
of the withdrawal threshold. It is the off-path interpretation of \(s_A\) once
the high type has crossed that threshold. Section~\ref{sec:belief-loop} closes
that belief loop.
